% =====================================================================
\chapter{Inversa, determinante y matrices ortogonales}
\label{tema:2}
% =====================================================================

Este tema estudia cuándo una transformación lineal puede deshacerse, cómo lo detecta el determinante y por qué las matrices ortogonales, cuya inversa es su traspuesta, ocupan un lugar privilegiado. Se sigue a Strang~\cite[§2.2, cap.~5 y §4.4]{strang-ila}, las clases~3, 17, 18 y~20 de 18.06~\cite{mit1806} y el §4.1 de~\cite{mml}.

\section{La matriz inversa}

\begin{definicion}[Inversa]
Una matriz cuadrada $A\in\R^{n\times n}$ es \emph{invertible} si existe $A^{-1}\in\R^{n\times n}$ tal que $AA^{-1}=A^{-1}A=I_n$. Una matriz no invertible se llama \emph{singular}.
\end{definicion}

\begin{proposicion}\label{prop:inversa-props}
Sean $A,B\in\R^{n\times n}$ invertibles.
\begin{enumerate}
  \item La inversa es única.
  \item $AB$ es invertible y $(AB)^{-1}=B^{-1}A^{-1}$.
  \item $A\T$ es invertible y $(A\T)^{-1}=(A^{-1})\T$.
\end{enumerate}
\end{proposicion}

\begin{proof}
(i) Si $B_1$ y $B_2$ son inversas de $A$, entonces $B_1=B_1(AB_2)=(B_1A)B_2=B_2$. (ii) $(AB)(B^{-1}A^{-1})=A(BB^{-1})A^{-1}=AA^{-1}=I$, y del mismo modo en el otro orden. (iii) Trasponiendo $AA^{-1}=I$ se obtiene $(A^{-1})\T A\T=I$, y análogamente $A\T(A^{-1})\T=I$.
\end{proof}

\begin{teorema}[Caracterizaciones de la invertibilidad]\label{teo:invertible}
Para $A\in\R^{n\times n}$ son equivalentes:
\begin{enumerate}
  \item $A$ es invertible;
  \item $\rango(A)=n$;
  \item $\Nuc(A)=\{0\}$, es decir, $Ax=0$ solo si $x=0$;
  \item $\det(A)\ne 0$;
  \item $0$ no es autovalor de $A$.
\end{enumerate}
\end{teorema}

\begin{proof}
(i)$\Rightarrow$(iii): si $Ax=0$, entonces $x=A^{-1}Ax=0$. (iii)$\Leftrightarrow$(ii): por el \cref{teo:rango-nulidad}, $\rango(A)=n-\dim\Nuc(A)$. (ii)$\Rightarrow$(i): si $\rango(A)=n$, las columnas de $A$ son una base de $\R^{n}$, luego cada $e_j$ se escribe $e_j=Ab_j$ de forma única; la matriz $B=[b_1\mid\cdots\mid b_n]$ cumple $AB=I$. Para ver que también $BA=I$, obsérvese que $A(BA-I)=(AB)A-A=0$, luego cada columna de $BA-I$ está en $\Nuc(A)=\{0\}$. La equivalencia con (iv) se justifica en la \cref{sec:determinante} y la equivalencia con (v) es inmediata a partir de la definición de autovalor del \cref{tema:3}: $0$ es autovalor si y solo si existe $x\ne 0$ con $Ax=0$.
\end{proof}

\begin{proposicion}[Inversa de orden 2]\label{prop:inversa2}
Sea $A=\begin{pmatrix}a&b\\c&d\end{pmatrix}$ con $ad-bc\ne 0$. Entonces
\[
A^{-1}=\frac{1}{ad-bc}\begin{pmatrix}d&-b\\-c&a\end{pmatrix}.
\]
\end{proposicion}

\begin{proof}
Basta multiplicar: $\begin{pmatrix}a&b\\c&d\end{pmatrix}\begin{pmatrix}d&-b\\-c&a\end{pmatrix}=\begin{pmatrix}ad-bc&0\\0&ad-bc\end{pmatrix}=(ad-bc)\,I$, y lo mismo en el orden inverso.
\end{proof}

\begin{ejemplo}\label{ej:inversa}
Para $A=\begin{pmatrix}2&1\\1&1\end{pmatrix}$, $ad-bc=2-1=1$ y $A^{-1}=\begin{pmatrix}1&-1\\-1&2\end{pmatrix}$. La primera fila de $AA^{-1}$ es $(2-1,\,-2+2)=(1,0)$ y la segunda, $(1-1,\,-1+2)=(0,1)$.
\end{ejemplo}

\begin{observacion}[Cálculo numérico]
Para resolver $Ax=b$ no se calcula $A^{-1}$: se factoriza $A$ (por ejemplo, $A=LU$) y se resuelven dos sistemas triangulares, lo que es más rápido y numéricamente más estable. En NumPy, \texttt{np.linalg.solve(A, b)} es preferible a \texttt{np.linalg.inv(A) @ b}. En los modelos de lenguaje las matrices de pesos no se invierten; lo que sí aparece constantemente es la ortogonalidad (\cref{sec:ortogonales}).
\end{observacion}

\section{El determinante}\label{sec:determinante}

Se adopta la presentación axiomática de~\cite[§5.1]{strang-ila} y de la clase~18 de~\cite{mit1806}.

\begin{definicion}[Determinante]\label{def:det}
El \emph{determinante} es la única función $\det\colon\R^{n\times n}\to\R$ que cumple:
\begin{enumerate}
  \item $\det(I_n)=1$;
  \item cambia de signo al intercambiar dos filas;
  \item es lineal en cada fila cuando las demás se mantienen fijas.
\end{enumerate}
\end{definicion}

De estas tres propiedades se deducen todas las demás~\cite[§5.1]{strang-ila}: si dos filas son iguales, el determinante es nulo; restar a una fila un múltiplo de otra no lo cambia; el determinante de una matriz triangular es el producto de su diagonal; $\det(A\T)=\det(A)$; y
\[
\det(AB)=\det(A)\det(B).
\]
La eliminación gaussiana lleva $A$ a una matriz triangular $U$ mediante operaciones que, a lo sumo, cambian el signo del determinante. Por tanto, $\det(A)\ne 0$ si y solo si todos los pivotes son no nulos, es decir, si y solo si $\rango(A)=n$. Esto completa la equivalencia (ii)$\Leftrightarrow$(iv) del \cref{teo:invertible}. Además, de $AA^{-1}=I$ se obtiene $\det(A^{-1})=1/\det(A)$.

\begin{proposicion}[Desarrollo por cofactores]\label{prop:laplace}
Para $A\in\R^{n\times n}$ y cualquier fila $i$,
\[
\det(A)=\sum_{j=1}^{n}(-1)^{i+j}a_{ij}\det(A_{ij}),
\]
donde $A_{ij}$ es la submatriz de orden $n-1$ obtenida al suprimir la fila $i$ y la columna $j$. En particular, $\det\begin{pmatrix}a&b\\c&d\end{pmatrix}=ad-bc$.
\end{proposicion}

Esta fórmula es el desarrollo de Laplace~\cite[teorema~4.2]{mml}; su demostración, a partir de la linealidad en la fila~$i$, está en~\cite[§5.2]{strang-ila}. Para $n=3$ equivale a la regla de Sarrus.

\begin{observacion}[Interpretación geométrica]
$\abs{\det(A)}$ es el factor por el que $A$ multiplica volúmenes: el cubo unidad se transforma en un paralelepípedo de volumen $\abs{\det(A)}$~\cite[§5.3]{strang-ila}. Un determinante nulo indica que el paralelepípedo es plano, es decir, que $A$ aplasta alguna dimensión. Para la matriz $M$ del \cref{ej:M}, $\det(M)=0$.
\end{observacion}

\section{Bases ortonormales y matrices ortogonales}\label{sec:ortogonales}

\begin{definicion}[Familia ortonormal]
Los vectores $q_1,\dots,q_k\in\R^{n}$ son \emph{ortonormales} si $q_i\T q_j=\delta_{ij}$, donde $\delta_{ij}$ vale~1 si $i=j$ y~0 en otro caso.
\end{definicion}

Una familia ortonormal es linealmente independiente: si $\sum_i c_iq_i=0$, multiplicando por $q_j\T$ se obtiene $c_j=0$.

\begin{proposicion}[Coordenadas en una base ortonormal]\label{prop:coord-ortonormal}
Si $\{q_1,\dots,q_n\}$ es una base ortonormal de $\R^{n}$, todo $x\in\R^{n}$ cumple
\[
x=\sum_{j=1}^{n}(q_j\T x)\,q_j\qquad\text{y}\qquad \norm{x}^{2}=\sum_{j=1}^{n}(q_j\T x)^{2}.
\]
Si la base es solo ortogonal (con $q_j\ne 0$), la coordenada $j$ es $c_j=q_j\T x/q_j\T q_j$.
\end{proposicion}

\begin{proof}
Escríbase $x=\sum_i c_iq_i$. Multiplicando por $q_j\T$, $q_j\T x=\sum_i c_i\,q_j\T q_i=c_j$. Además, $\norm{x}^{2}=x\T x=\sum_{i,j}c_ic_j\,q_i\T q_j=\sum_j c_j^{2}$. En el caso ortogonal, $q_j\T x=c_j\,q_j\T q_j$.
\end{proof}

La segunda igualdad es el teorema de Pitágoras en $\R^{n}$. Calcular coordenadas en una base ortonormal no requiere resolver ningún sistema: basta un producto escalar por coordenada.

\begin{definicion}[Matriz ortogonal]
Una matriz $Q\in\R^{n\times n}$ es \emph{ortogonal} si sus columnas forman una base ortonormal de $\R^{n}$, es decir, si $Q\T Q=I_n$. En ese caso, $Q^{-1}=Q\T$.
\end{definicion}

\begin{ejemplo}[Giros del plano]\label{ej:giro}
Para cada ángulo $\theta$, la matriz
\[
Q_\theta=\begin{pmatrix}\cos\theta&-\sin\theta\\ \sin\theta&\cos\theta\end{pmatrix}
\]
es ortogonal: sus columnas tienen norma $\cos^{2}\theta+\sin^{2}\theta=1$ y su producto escalar es $-\cos\theta\sin\theta+\sin\theta\cos\theta=0$. Su inversa, $Q_\theta\T=Q_{-\theta}$, es el giro en sentido contrario, y $\det(Q_\theta)=1$. La reflexión $\begin{pmatrix}1&0\\0&-1\end{pmatrix}$ también es ortogonal, con determinante $-1$.
\end{ejemplo}

\begin{proposicion}\label{prop:ortogonal-isometria}
Si $Q$ es ortogonal, entonces $(Qx)\T(Qy)=x\T y$ y $\norm{Qx}=\norm{x}$ para todos $x,y$, y $\det(Q)=\pm 1$.
\end{proposicion}

La demostración se propone en el \cref{eje:2-isometria}. Las matrices ortogonales son, por tanto, las transformaciones lineales que conservan longitudes y ángulos: giros y reflexiones.

\begin{observacion}[Gram--Schmidt]
A partir de cualquier base $a_1,\dots,a_n$ se obtiene una base ortonormal del mismo espacio restando a cada $a_k$ su proyección sobre los vectores ya construidos y normalizando: $w_k=a_k-\sum_{i<k}(q_i\T a_k)\,q_i$ y $q_k=w_k/\norm{w_k}$~\cite[§4.4]{strang-ila}. En forma matricial, el proceso da la factorización $A=QR$ con $R$ triangular superior.
\end{observacion}

\begin{observacion}[Ortogonalidad en un LLM]
Las matrices ortogonales aparecen en varios puntos de un Transformer. La codificación posicional rotatoria (RoPE, módulo~14) aplica a cada par de coordenadas de las consultas y las claves un giro $Q_\theta$ cuyo ángulo depende de la posición; como los giros conservan la norma, solo cambian los ángulos entre vectores. Las matrices $U$ y $V$ de la SVD (\cref{tema:4}) son ortogonales, y algunas inicializaciones de pesos usan matrices ortogonales aleatorias porque conservan la norma de las activaciones.
\end{observacion}

\begin{lecturas}
\item Strang~\cite{strang-ila}, §2.2 «Elimination Matrices and Inverse Matrices», cap.~5 «Determinants» y §4.4 «Orthonormal Bases and Gram-Schmidt».
\item MIT 18.06~\cite{mit1806}, clases~3 «Multiplication and inverse matrices», 17 «Orthogonal matrices and Gram-Schmidt», 18 «Properties of determinants» y 20 «Cramer's rule, inverse matrix, and volume».
\item Deisenroth, Faisal y Ong~\cite{mml}, §4.1 «Determinant and Trace».
\item Goodfellow, Bengio y Courville~\cite{gbc}, §2.3 «Identity and Inverse Matrices» y §2.11 «The Determinant».
\end{lecturas}

\section*{Ejercicios del tema 2}
\addcontentsline{toc}{section}{Ejercicios del tema 2}

\begin{ejercicio}{1}{eje:2-equiv}
Indique, sin calcular inversas, cuáles de las siguientes matrices son invertibles y qué condición del \cref{teo:invertible} lo justifica: (a) la matriz $M$ del \cref{ej:M}; (b) $Q_\theta$ del \cref{ej:giro}; (c) $\begin{pmatrix}1&2\\2&4\end{pmatrix}$; (d) una matriz $A\in\R^{4\times 4}$ que cumple $A(1,0,-1,2)=0$; (e) $\diag(3,\,0{,}5,\,-1)$.
\end{ejercicio}

\begin{ejercicio}{2}{eje:2-inversa}
Calcule la inversa de $A=\begin{pmatrix}1&2\\3&4\end{pmatrix}$, compruebe que $AA^{-1}=I$ y verifique que $\det(A^{-1})=1/\det(A)$.
\end{ejercicio}

\begin{ejercicio}{2}{eje:2-mml}
(a) Calcule el determinante de
\[
A=\begin{pmatrix}1&3&5\\2&4&6\\0&2&4\end{pmatrix}
\]
mediante el desarrollo por la primera fila y mediante la regla de Sarrus. (b) Determine, si existe, la inversa de $B=\begin{pmatrix}2&3&4\\3&4&5\\4&5&6\end{pmatrix}$. En caso de que no exista, dé un vector no nulo de $\Nuc(B)$.
\fuente{adaptado de Deisenroth, Faisal y Ong~\cite{mml}, ejercicios~4.1 y~2.8(a).}
\end{ejercicio}

\begin{ejercicio}{2}{eje:2-giros}
(a) Compruebe que $Q_\alpha Q_\beta=Q_{\alpha+\beta}$. (b) Calcule $Q_{\pi/6}\,(2,0)$ y compruebe que su norma es~2. (c) Explique por qué el producto escalar $(Q_{\alpha}x)\T(Q_{\beta}y)$ depende de $\alpha$ y $\beta$ solo a través de $\beta-\alpha$. Esta es la propiedad que aprovecha RoPE para que la atención dependa de la posición relativa.
\end{ejercicio}

\begin{ejercicio}{3}{eje:2-haar}
En $\R^{4}$, la base de Haar está formada por $h_1=(1,1,1,1)$, $h_2=(1,1,-1,-1)$, $h_3=(1,-1,0,0)$ y $h_4=(0,0,1,-1)$. (a) Compruebe que es ortogonal y normalícela. (b) Escriba la señal $x=(4,2,5,5)$ en esa base. (c) Comprima $x$ conservando solo los dos primeros coeficientes y calcule el error cuadrático. Justifique, mediante la \cref{prop:coord-ortonormal}, que el error es la suma de los cuadrados de los coeficientes descartados multiplicados por $\norm{h_j}^{2}$.
\fuente{adaptado de MIT 18.06SC~\cite{mit1806sc}, sesión «Change of Basis; Image Compression», problema~31.1 (base de 8 vectores reducida a 4; apartados (b) y (c) añadidos).}
\end{ejercicio}

\begin{ejercicio}{3}{eje:2-isometria}
Demuestre la \cref{prop:ortogonal-isometria}. Demuestre además que el producto de dos matrices ortogonales es ortogonal.
\end{ejercicio}
